\paragraph*{Bad vertices with respect to the ordered \toughpair\ $(A,B)$.}
Given an ordered \toughpair\ $(A,B)$ in a digraph $D$,
we say that a vertex $v \in V(D) \setminus V(A\cup B)$ is {\em bad} with respect to $(A,B)$ if $v \in N^+_{D^{\star}}(\tail(A)) \cap N^-_{D^{\star}}(\head(A)) \cap N^+_{D^{\star}}(\tail(B)) \cap N^-_{D^{\star}}(\head(B))$. 

\begin{lemma}[Eliminating bad vertices]\label{lem:badvertices}
Let $D$ be a
digraph such that $D$ contains an ordered $t^2$-\toughpair. Then $D$ contains an ordered $t$-\toughpair\ $(A',B')$ such that there is no bad vertex in $V(D) \setminus V(A' \cup B')$ with respect to $(A',B')$.
\end{lemma}

\begin{proof}
Let $(A,B)$ be an ordered $t^2$-\toughpair\ in $D$ and
let $t'=t^2$.
Let $A=((w_1,x_1), \ldots, (w_{t'},x_{t'}))$ and $B=((y_1,z_1), \ldots, (y_{t'},z_{t'}))$.
For each vertex $v \in V(H) \setminus V(A \cup B)$ which is bad with respect to $(A,B)$, let $W(v)$ be the largest index in $[t']$ such that $(a_{W(v)},v) \in E(D)$, $X(v)$ be the smallest index in $[t']$ such that $(v,b_{X(v)}) \in E(D)$, $Y(v)$ be the largest index in $[t']$ such that $(c_{Y(v)},v) \in E(D)$, and $Z(v)$ be the smallest index in $[t']$ such that $(v,d_{Z(v)}) \in E(D)$.
For every $p,q \in [\sqrt{t'}]$, let $V_{p,q}$ be the set of bad vertices $v$ with respect to $(A,B)$ such that
$W(v),X(v) \in \{(p-1)\sqrt{t'}+1, \ldots, p\sqrt{t'}\}$, and $Y(v),Z(v) \in \{(q-1)\sqrt{t'}, \ldots, q\sqrt{t'}\}$.

If there exists $p,q \in [\sqrt{t'}]$ such that $V_{pq}=\emptyset$, then let $A'=((w_{(p-1)\sqrt{t'}+1},x_{(p-1)\sqrt{t'}}+1), \ldots, (x_{p\sqrt{t'}}, x_{p\sqrt{t'}}))$ and 
$B'=((c_{(q-1)\sqrt{t'}+1},z_{(q-1)\sqrt{t'}+1}),\ldots, (c_{q\sqrt{t'}},z_{q\sqrt{t'}}))$. Then from the definition of $V_{p,q}$, $(A',B')$ is an ordered $\sqrt{t'}$-\toughpair\ such that there are no bad vertices with respect to it.

Without loss of generality, assume that for each $p,q \in [\sqrt{t'}]$, $V_{pq}\neq \emptyset$. Let $v_{p,q}$ denote an arbitrarily fixed vertex in $V_{pq}$. 
Let $A'=((w_{W(v_{2,1}}, v_{2,1}), \ldots, (w_{W(v_{\sqrt{t'},1})}, v_{\sqrt{t'},1}))$ and
 $B'=((y_{Y(v_{1,2})}, v_{1,2}),\ldots, (y_{Y(v_{\sqrt{t'},1})},v_{1,\sqrt{t'}} ))$. We will now show that $(A',B')$ is an ordered \toughpair\ in $D$. Moreover, we will show that there is no $(\tail(B'),\head(A'))$-path in $D$, thereby concluding that there is no bad vertex with respect to $(A',B')$. Note that the edges of $A',B'$ are minimal edges because they belong to $D$ which is a reachability minimal digraph.

\paragraph*{Weak independence of $A'$ (resp.~$B'$).}  We show that  there is no $(\head(A'),\tail(A'))$ path (resp.~no $(\head(B'),\tail(B'))$ path) in $D$. Recall that every vertex in $\head(A')$ is a bad vertex with respect to $(A,B)$.
Thus, each vertex of $\tail(A')$ is an out-neighbour of some vertex of $\tail(B')$. 
Thus, a $(\head(A'),\tail(A'))$-path in $D$ implies a $(\tail(B),\tail(A'))$-path in $D$. Since $\tail(A') \subseteq \tail(A)$, this contradicts that every edge of $A$ is strongly independent with every edge of $B$. 
The other case is symmetric.

\paragraph*{Weak independence between $A'$ and $B'$.} We show that a pair of edges containing an edge of $A'$ and an edge of $B'$ is weakly independent.
To see that there is no $(\head(B'),\tail(A'))$-path, 
recall that every vertex of $\head(B')$ is a bad vertex with respect to $(A,B)$.
Thus, each vertex of $\head(B')$ is an out-neighbour of some vertex of $\tail(B)$. Thus, a $(\head(B'),\tail(A'))$-path in $D$ implies a $(\tail(B),\head(A'))$-path in $D$. Since $\tail(A') \subseteq \tail(A)$, this is a contradiction. One can symmetrically prove that there is no
$(\head(A'),\tail(B'))$-path.

\paragraph*{Strong independence of $A'$ and $B'$.} We prove the following three statements.

{\it There is no $(\tail(A'),\tail(B'))$-path and no $(\tail(B'),\tail(A'))$-path.} Since $\tail(A') \subseteq \tail(A)$ and $\tail(B') \subseteq \tail(B)$, and $(A,B)$ is a  \toughpair, we conclude that there is no $(\tail(A'),\allowbreak
\tail(B'))$-path and no $(\tail(B'),\tail(A'))$-path.

{\it There is no $(\head(A'),\head(B'))$-path.} For the sake of contradiction, suppose there exists $v_{i,1} \in \head(B')$ and $v_{1,j} \in \head(B')$, $i,j \neq 1$, such that there exists a $(v_{i,1}, v_{1,j})$-path in $D$. By definition $w(v_{i,1}) \in \{2\sqrt{t'}+1, \ldots, t'\}$ and $X(v_{1,j}) \in \{1, \ldots, \sqrt{t'}\}$. Also, $(w_{W(v_{i,1})}, v_{i,1}) \in E(D)$ and $(v_{1,j}, x_{X(v_{1,j})}) \in E(D)$. Thus, a $(v_{i,1}, v_{1,j})$-path in $D$ implies a $(w_{W(v_{i,1})}, y_{Y(v_{1,j})})$-path in $D$. Since $X(v_{1,j}) < W(v_{i,1})$, this is a contradiction.

{\it There is no $(\head(B'),\head(A'))$-path.} For the sake of contradiction, suppose there exists $v_{i,1} \in \head(A')$ and $v_{1,j} \in \head(B')$, $i,j \neq 1$, such that there exists a $(v_{1,j}, v_{i,1})$-path in $D$. By definition $Y(v_{1,j}) \in \{2\sqrt{t'}+1, \ldots, t'\}$ and $Z(v_{i,1}) \in \{1, \ldots, \sqrt{t'}\}$. Also, $(y_{Y(v_{1,j})}, v_{1,j}) \in E(D)$ and $(v_{i,1}, z_{Z(v_{i,1})}) \in E(D)$. Thus, a $(v_{1,j}, v_{i,1})$-path in $D$ implies a $(y_{Y(v_{1,j})}, z_{Z(v_{i,1})})$-path in $D$. Since $Z(v_{i,1}) < Y(v_{1,j})$, this is a contradiction.

\paragraph*{Ordered condition on $A'$ (resp.~$B'$).}
Consider vertices $v_{i,1}$ and $v_{j,1}$ such that $i \leq j$. Suppose for the sake of contradiction that there exists a $(v_{j,1},v_{i,1})$-path in $D$. Recall $(w_{W(v_{j,1})}, v_{j,1}) \in E(D)$ and $(v_{i,1},x_{X(v_{i,1})}) \in E(D)$. Thus, there exists a $(w_{W(v_{j,1})}, y_{Y(v_{i,1})})$-path in $D$. Since $i \leq j$, $W(v_{j,1}) \leq X(v_{i,1})$. This is a contradiction. The other case can be proved symmetrically.

\paragraph*{There is no $(\tail(B'),\head(A'))$-path.} For the sake of contradiction, let $v_{i,1} \in \head(A')$, $i\neq 1$, and $y_j \in \tail(B')$ such that there exists a $(y_j, v_{i,1})$-path in $D$. Since $y_j \in \tail(B')$, $j \in \{2\sqrt{t'} +1, \ldots, t'\}$. Also $Z(v_{i,1}) \in \{1, \ldots, \sqrt{t}\}$. Since $(v_{i,1}, z_{Z(v_{i,1})}) \in E(D)$, we conclude that there is a $(y_j, z_{Z(v_{i,1})})$-path. This is a contradiction which concludes the proof.
\end{proof}